In my weld-seam HMI, modified-DH kinematics connected joint commands to an OpenGL digital twin. My later CR5 work added a usable fixed-view calibration with documented accuracy limits, workspace checks, and blocking short-horizon ServoP execution. These projects gave me experience joining geometry and control, but they also exposed my next gap: integrating automation, sensing, and modern robotics into a disciplined process for physical validation. The current Robotics stream in CityUHK's MSc in Mechanical Engineering addresses that progression.

If that stream structure and its courses continue into my intake, MNE6007 Advanced Automation Technology would deepen the automation logic behind integrated machines, MNE6126 Sensors for Robotics, AI, and Control Systems would strengthen the link from measured state to control decisions, and MNE6130 Modern Robotics would place both within contemporary robotic systems. If I selected the alternative dissertation route and received approval, the optional nine-credit MNE6008 Dissertation could support a bounded study of how calibration limits and execution timing affect repeatable robot behavior. I would bring coordinate-modeling and real-hardware integration experience to that work, then apply stronger physical-validation methods in robotics or intelligent-machine R\&D.
